%=========================== APÊNDICES ===========================
\newpage
\section*{Apêndice A --- Termo de Consentimento Livre e Esclarecido}
\addcontentsline{toc}{section}{APÊNDICE A --- TCLE}

% RASCUNHO. Conferir com o orientador se a instituição exige submissão
% ao Comitê de Ética em Pesquisa (CEP) e se há modelo próprio de TCLE.
% Este texto é também o que a tela de cadastro da aplicação apresenta.

\begingroup\singlespacing\small

\noindent\textbf{Título do estudo:} \tcctitulo

\noindent\textbf{Pesquisadora:} \tccautor \quad \textbf{Orientador:} \tccorientador

\vspace{0.4cm}
\noindent\textbf{Do que se trata.} Você está sendo convidado(a) a participar de
uma pesquisa de Trabalho de Conclusão de Curso que avalia um protótipo de
Realidade Virtual. O sistema apresenta animais comumente associados a fobias
--- cobra, rato, barata, aranha e sapo --- em ambientes virtuais, organizados em
níveis crescentes de aproximação.

\vspace{0.3cm}
\noindent\textbf{O que não é.} Este sistema \textbf{não é tratamento
psicológico ou médico}, \textbf{não realiza diagnóstico} e \textbf{não
substitui acompanhamento profissional}. A avaliação restringe-se a usabilidade,
imersão, interação e percepção.

\vspace{0.3cm}
\noindent\textbf{O que você fará.} Preencherá um cadastro e um questionário
curto, usará a aplicação pelo tempo que desejar e responderá a um questionário
final. A sessão dura cerca de [N] minutos.

\vspace{0.3cm}
\noindent\textbf{Dados coletados.} E-mail, sexo e data de nascimento; o grau de
medo que você declara para cada animal; e medidas automáticas de uso, como
distância mantida até o estímulo, movimentação e níveis percorridos.

\vspace{0.3cm}
\noindent\textbf{Como seus dados são tratados.} Seu e-mail serve apenas para
identificar o seu cadastro e \textbf{não acompanha os dados de análise}: nos
resultados você é identificado(a) por um código (P01, P02, \ldots). Nenhum dado
que permita identificá-lo(a) será publicado.

\vspace{0.3cm}
\noindent\textbf{Riscos e desconforto.} O contato com imagens de animais pode
causar desconforto. \textbf{Você pode interromper a qualquer momento}, sem
precisar justificar e sem prejuízo de qualquer natureza. O uso de óculos de
Realidade Virtual pode provocar tontura passageira.

\vspace{0.3cm}
\noindent\textbf{Benefícios.} Não há benefício direto. A participação contribui
para a avaliação de uma ferramenta de pesquisa.

\vspace{0.3cm}
\noindent\textbf{Voluntariedade.} A participação é voluntária. Você pode
desistir a qualquer momento e solicitar a exclusão dos seus dados pelo contato
[E-MAIL DA PESQUISADORA].

\vspace{0.6cm}
\noindent Declaro que li e compreendi as informações acima e concordo em
participar.

\vspace{0.8cm}
\noindent\rule{8cm}{0.4pt}\hfill\rule{4cm}{0.4pt}\\
Assinatura do(a) participante \hfill Data

\endgroup

\newpage
\section*{Apêndice B --- Questionário inicial de perfil}
\addcontentsline{toc}{section}{APÊNDICE B --- Questionário inicial}

% Aplicado no cadastro, dentro da própria aplicação.
% Curto de propósito: a pessoa está curiosa para entrar.

\begin{enumerate}
  \item Sexo: ( ) Feminino \quad ( ) Masculino \quad ( ) Outro \quad ( ) Prefiro não informar
  \item Data de nascimento: \rule{3cm}{0.4pt}
  \item Curso e período: \rule{6cm}{0.4pt}
  \item Já usou óculos de Realidade Virtual? ( ) Sim \quad ( ) Não
  \item De 0 a 10, quanto medo você sente de cada animal?
  \begin{itemize}
    \item Cobra: \rule{1.5cm}{0.4pt} \quad Rato: \rule{1.5cm}{0.4pt} \quad
          Barata: \rule{1.5cm}{0.4pt} \quad Aranha: \rule{1.5cm}{0.4pt} \quad
          Sapo: \rule{1.5cm}{0.4pt}
  \end{itemize}
\end{enumerate}

\newpage
\section*{Apêndice C --- Questionário final de percepção}
\addcontentsline{toc}{section}{APÊNDICE C --- Questionário final}

% Respostas de 0 a 10, salvo indicação em contrário.

\begin{enumerate}
  \item O quanto você se sentiu presente dentro do ambiente virtual?
  \item O quanto o animal pareceu real para você?
  \item O quanto foi fácil aprender a usar o sistema?
  \item O quanto foi fácil se deslocar pelo ambiente?
  \item As instruções estavam claras?
  \item Você sentiu desconforto físico (tontura, enjoo)?
  \item O quanto você se sentiu no controle do próprio ritmo?
  \item Em que nível você sentiu maior desconforto? ( ) 1 ( ) 2 ( ) 3 ( ) 4 ( ) 5 ( ) nenhum
  \item Você interromperia a experiência em algum momento? Por quê?
  \item O que você mudaria no sistema?
\end{enumerate}
